\documentclass[12pt]{article}
\usepackage[utf8]{inputenc}
\usepackage[T1]{fontenc}
\usepackage{hyperref}
\title{Vasavi Language Resources for K-12 Curriculum Development}
\author{Compiled by Claude Sonnet 5 (claude-sonnet-5) for Andr\'as Kornai}
\date{}
\begin{document}
\maketitle

\section{Introduction}
Vasavi (ISO 639-3 \texttt{vas}, Glottocode vasa1239) is a Western Indo-Aryan (Bhil) language of the Vasava/Vasave people, spoken mainly in Gujarat (Bharuch district) and Maharashtra (Nandurbar and Dhule districts). It is distinct from Bhili. Speaker counts in sources differ widely (Wikipedia cites 187,000 in 2011; Global Recordings lists about 900,000).

\section{Primary Educational Resources}
No primers, textbooks or school programmes were found in the search results.

\section{Linguistic and Reference Materials}
\begin{itemize}
\item Wikipedia overview, including dialects (e.g.\ Dehvaliya): \url{https://en.wikipedia.org/wiki/Vasavi_language}
\item Glottolog language page: \url{https://glottolog.org/resource/languoid/id/vasa1239.bigmap.html}
\end{itemize}

\section{Religious and Community Resources}
\begin{itemize}
\item Scripture Earth listing for Vasavi: \url{https://www.scriptureearth.org/00eng.php?iso=vas}
\item Joshua Project profile of Vasave Bhil (population by district): \url{https://joshuaproject.net/people_groups/21022}
\item Wikipedia on the Vasava people: \url{https://en.wikipedia.org/wiki/Vasava_people}
\end{itemize}

\section{Digital and Audio Resources}
\begin{itemize}
\item Global Recordings Network page for Vasavi: \url{https://globalrecordings.net/en/language/vas}
\end{itemize}

\section{Curriculum Development Context}
The language is reported as not taught in schools. The search found no dictionary, grammar or teaching material; the listings above are general references and missionary/audio materials whose contents were not inspected. Any curriculum would require new material, and the Vasavi/Bhili distinction should be kept in view.
\end{document}